\chapter{Inspecting and Changing Data}\label{ch:gdb-data}

You are stopped in the right frame. Now the actual work: finding out what the
program thinks is true.

\section{\texorpdfstring{\cmd{print} and Format Letters}{print and Format Letters}}\label{sec:gdb-print}

\cmd{print} evaluates an expression in the language of the current frame ---
real C, with your program's variables, operators and functions.

\begin{keys}[0.26]
	\cmd{print n}              & A variable                                        \\
	\cmd{print a[i]}           & An array element                                  \\
	\cmd{print c->next->name}  & A chain of dereferences                           \\
	\cmd{print n * 2 + 1}      & Arithmetic                                        \\
	\cmd{print s == NULL}      & A comparison: prints 1 or 0                       \\
	\cmd{print sizeof(struct node)} & Ask the type system                          \\
	\cmd{print func(3)}        & Call it (\autoref{sec:gdb-force})                 \\
\end{keys}

A slash and a letter changes the representation:

\begin{keys}[0.16]
	\cmd{/x} & Hexadecimal                        \\
	\cmd{/d} & Signed decimal                     \\
	\cmd{/u} & Unsigned decimal                   \\
	\cmd{/o} & Octal                              \\
	\cmd{/t} & Binary --- for flags and masks     \\
	\cmd{/c} & As a character                     \\
	\cmd{/f} & As a float                         \\
	\cmd{/s} & As a string                        \\
	\cmd{/a} & As an address, with the nearest symbol \\
\end{keys}

\begin{gdbcode}
(gdb) print flags
$1 = 37
(gdb) print/x flags
$2 = 0x25
(gdb) print/t flags
$3 = 100101
(gdb) print/a main
$4 = 0x1149 <main>
\end{gdbcode}

\begin{note}
	\cmd{print/t} is the one to remember. Bit flags read as decimal are
	meaningless; in binary the answer is usually visible at a glance.
\end{note}

Several settings make output readable, and belong in your \file{.gdbinit}:

\begin{keys}[0.30]
	\cmd{set print pretty on}      & One struct member per line, indented         \\
	\cmd{set print array on}       & Arrays one element per line                  \\
	\cmd{set print elements 0}     & Do not truncate long arrays or strings       \\
	\cmd{set print null-stop on}   & Stop printing a char array at the first NUL  \\
	\cmd{set print union off}      & Do not expand unions inside structs          \\
	\cmd{set pagination off}       & Do not stop every screenful                  \\
\end{keys}

\begin{gdbcode}
(gdb) set print pretty on
(gdb) print *cfg
$5 = {
  name = 0x555555556004 "server",
  port = 8080,
  timeout = 30,
  next = 0x0
}
\end{gdbcode}

\section{\texorpdfstring{\cmd{x}: Examining Memory}{x: Examining Memory}}\label{sec:gdb-x}

\cmd{print} shows a \emph{value}; \cmd{x} shows the \emph{bytes at an address},
which is what you need when the value is not what its type claims.

\begin{center}
	\ttfamily\large
	x/\{count\}\{format\}\{size\} \{address\}
\end{center}

\begin{keys}[0.22]
	\cmd{x/16xb buf}   & 16 bytes, in hex                                  \\
	\cmd{x/8xw buf}    & 8 words (4 bytes each), in hex                    \\
	\cmd{x/4xg \$sp}   & 4 giants (8 bytes each) from the stack pointer    \\
	\cmd{x/s s}        & As a NUL-terminated string                        \\
	\cmd{x/10c buf}    & 10 characters                                     \\
	\cmd{x/5i \$pc}    & 5 instructions                                    \\
	\cmd{x/20dw arr}   & 20 words as signed decimal                        \\
\end{keys}

Sizes are \cmd{b}yte, \cmd{h}alfword (2), \cmd{w}ord (4) and \cmd{g}iant (8).

\begin{gdbcode}
(gdb) x/16xb buf
0x7fffffffe2a0: 0x68 0x65 0x6c 0x6c 0x6f 0x00 0xff 0xff
0x7fffffffe2a8: 0xff 0xff 0xff 0xff 0x00 0x00 0x00 0x00
(gdb) x/s buf
0x7fffffffe2a0: "hello"
\end{gdbcode}

\begin{note}
	\cmd{x} remembers its format, so a bare \cmd{x} after the above shows the next
	16 bytes in the same layout --- and \key{<CR>} repeats it. Walking a buffer is
	one command and then the Enter key.
\end{note}

\begin{moral}
	Reach for \cmd{x} when you distrust the type. A \texttt{char *} that prints as
	garbage, a struct whose fields look implausible, a buffer you suspect was
	overrun --- \cmd{x/32xb} shows what is really there, with no interpretation
	layered on top.
\end{moral}

\section{Types}\label{sec:gdb-types}

\begin{keys}[0.26]
	\cmd{ptype cfg}          & The full definition of the variable's type        \\
	\cmd{ptype struct node}  & \dots{} of a named type                           \\
	\cmd{whatis cfg}         & Just the type name, without expanding it          \\
	\cmd{ptype/o struct node} & With byte offsets and sizes --- shows the padding \\
	\cmd{print sizeof(cfg)}  & Its size                                          \\
	\cmd{info types node}    & Which types match a pattern                       \\
\end{keys}

\begin{gdbcode}
(gdb) ptype/o struct node
/* offset      |    size */  type = struct node {
/*      0      |       8 */    char *name;
/*      8      |       4 */    int port;
/* XXX  4-byte hole      */
/*     16      |       8 */    struct node *next;
                               /* total size (bytes):   24 */
                             }
\end{gdbcode}

\begin{note}
	\cmd{ptype/o} is how you see padding, which matters when a struct is written to
	a file or sent over a socket. The four-byte hole above is why
	\cmd{sizeof(struct node)} is 24 and not 20 --- and why writing the struct
	verbatim produces a format that differs between compilers.
\end{note}

\section{\texorpdfstring{\cmd{display}}{display}}\label{sec:gdb-display}

\begin{keys}[0.26]
	\cmd{display i}        & Print \texttt{i} automatically at every stop       \\
	\cmd{display/x flags}  & \dots{} with a format                              \\
	\cmd{display/i \$pc}   & The next instruction, every time                   \\
	\cmd{info display}     & What is being displayed                            \\
	\cmd{undisplay 1}      & Stop displaying it                                 \\
\end{keys}

\begin{gdbcode}
(gdb) display i
1: i = 0
(gdb) display total
2: total = 0
(gdb) next
15              total += weight[i];
1: i = 0
2: total = 0
(gdb) next
14          for (int i = 0; i < n; i++) {
1: i = 0
2: total = 12
\end{gdbcode}

\begin{moral}
	Two or three \cmd{display}s turn stepping into a live view of the variables you
	care about. Set them once, then hold \key{<CR>} and watch --- it is the closest
	GDB gets to a graphical debugger's watch pane, and it is faster to set up.
\end{moral}

\section{Changing Variables and Registers}\label{sec:gdb-set}

\begin{keys}[0.28]
	\cmd{set var n = 42}       & Assign to a variable                             \\
	\cmd{set var p = 0}        & Make a pointer null, to test the error path      \\
	\cmd{set var s[0] = 'X'}   & Modify one element                               \\
	\cmd{set var \$rax = 1}    & Assign to a register                             \\
	\cmd{print n = 42}         & Assigns too, and shows the result                \\
\end{keys}

\begin{note}
	Use \cmd{set var}, not bare \cmd{set}. \cmd{set width = 80} is GDB's own
	\cmd{set width} setting, not an assignment to your variable named
	\texttt{width}; \cmd{set var width = 80} is unambiguous. This bites people
	roughly once, memorably.
\end{note}

\begin{eg}
	Testing an error path that is hard to trigger for real:
\begin{gdbcode}
(gdb) break write_output
(gdb) run
Breakpoint 1, write_output (f=0x5555555592a0, n=1024) at io.c:55
(gdb) set var f = 0            # pretend the file failed to open
(gdb) continue
error: cannot write output     # good -- the check works
\end{gdbcode}
\end{eg}

\section{Pointers, Arrays and Structs}\label{sec:gdb-agg}

\begin{keys}[0.28]
	\cmd{print *p}           & Dereference                                        \\
	\cmd{print p->field}     & A member through a pointer                         \\
	\cmd{print *arr@10}      & Ten elements starting at \texttt{arr}              \\
	\cmd{print *p@n}         & \texttt{n} elements --- \texttt{n} may be a variable \\
	\cmd{print arr[3]@4}     & Elements 3 to 6                                    \\
	\cmd{print (struct node *)ptr} & Cast, when the type was lost                 \\
	\cmd{print \{int\}0x7fff2a0} & Read that address as an \texttt{int}           \\
	\cmd{print *node->next->next} & Follow a list by hand                         \\
\end{keys}

\begin{definition}[The \texttt{@} operator]
	\texttt{expr@n} means ``treat this as the first of \texttt{n} consecutive
	elements''. It exists because a pointer in C carries no length, so GDB cannot
	know how much of a dynamically allocated array to print.
\end{definition}

\begin{gdbcode}
(gdb) print arr
$1 = (int *) 0x5555555592a0        # just an address: no length information
(gdb) print *arr@5
$2 = {3, 1, 4, 1, 5}
(gdb) print *arr@n                 # n is the program's own length variable
$3 = {3, 1, 4, 1, 5, 9, 2, 6}
\end{gdbcode}

\begin{moral}
	\texttt{*p@n} is the command that makes GDB usable on C. Without it every
	heap array prints as a bare address, and half the state of your program is
	invisible.
\end{moral}

\section{Convenience Variables and Value History}\label{sec:gdb-convvar}

\begin{keys}[0.26]
	\cmd{\$1}, \cmd{\$2}   & Earlier results. Every \cmd{print} numbers its output  \\
	\cmd{\$}               & The most recent value                                  \\
	\cmd{\$\$}             & The one before that                                    \\
	\cmd{\$\$3}            & Three back                                             \\
	\cmd{set \$i = 0}      & A convenience variable --- yours, not the program's     \\
	\cmd{set \$p = head}   & Useful for walking a data structure                    \\
	\cmd{show values}      & The last ten history entries                           \\
\end{keys}

\begin{eg}[Walking a linked list]
\begin{gdbcode}
(gdb) set $n = head
(gdb) print $n->name
$1 = 0x555555556004 "first"
(gdb) set $n = $n->next
(gdb) print $n->name
$2 = 0x55555555600a "second"
\end{gdbcode}
	Now the last two lines repeat with \key{<CR>} --- two keystrokes per node. For
	a longer list, a loop does it in one command:
\begin{gdbcode}
(gdb) set $n = head
(gdb) while $n != 0
 >printf "%s -> ", $n->name
 >set $n = $n->next
 >end
first -> second -> third ->
\end{gdbcode}
\end{eg}

\begin{note}
	Convenience variables are typeless and persist across \cmd{run}, which makes
	them the right place to stash an address you want to keep watching across
	several runs of the program. They never collide with your program's own names,
	because your program cannot have a variable beginning with \texttt{\$}.
\end{note}
